\begin{multicols}{2}
\scieloSectionContainer{INTRODUÇÃO}
\par{}As leucemias mieloides agudas (LMA) constituem um
grupo heterogêneo de doenças caracterizado pela expansão
clonal de progenitores hematopoéticos imaturos na medula
óssea levando ao bloqueio da hematopoese normal 1.
A leucemia promielocítica aguda (LPA) é um subtipo da
LMA, caracterizada pelo predomínio de células com morfologia
característica pela translocação t(15,17) e por coagulopatia,
que combina coagulação intravascular disseminada e
fibrinólise 2.

\scieloSectionContainer{RELATO DO CASO}

\par{}Escolar, 8 anos e 2 meses, com diagnóstico pregresso
de leucemia promielocítica aguda com tratamento completo
concluído há 1 ano e sem intercorrências durante este período.
Deu entrada no PSI do HC-UFTM, no dia 12 de abril de
2014, com queixa de sangramento gengival, equimoses em
membros inferiores e plaquetopenia. Em hemograma inicial,
visualizados 70% de blastos, levantada hipótese diagnóstica de
recidiva da doença, sendo então iniciado ATRA 25 mg/m 2 /dia
e suporte hemoterápico. Mielograma evidenciou medula
óssea hiperplásica e homogênea, com 90% de mieloblastos e
imunofenotipagem compatível com leucemia promielocítica
aguda, confirmada com identificação do gene de fusão
PML-RARA

\lipsum

\scieloSectionContainer{Materials and Methods}

\scieloSubSectionContainer{Animals}

\par{}Specimens representing 12 species of polyclads were
collected from different sites in the Caribbean. In addition,
a deep-sea specimen of \textit{Anocellidus profundus} Quiroga,
Bolaños and Litvaitis, 2006, from the North Pacific was
also included in the analysis (specimen courtesy of Dr.
Janet Voight, Field Museum of Natural History, Chicago,
Illinois, USA). For details of collection, fixation, and spe-
cies identification see Litvaitis \textit{et al}. (2010).

\lipsum
\scieloSectionContainer{Considerações finais}

\par{}O processo de metropolização tem sido acompanhado
de um importante intensificação das desigualdades sociais, potencializando a difusão de práticas securitárias.
A partir da percepção de um aumento da violência nas
grandes metrópoles, a segurança é uma mercadoria
capaz de ser adquirida e consumida. Nas periferias
metropolitanas, torna-se cada vez mais comum a
disseminação de empreendimentos exclusivos cujo
acesso é controlado através de tecnologias de segurança


\lipsum

\scieloImageContainerOneCol{0.5\textwidth}{img1.jpg}{meu titulo}

%\begin{center}
%{
%\vspace{8mm}
%\centerline{
%\includegraphics[width=\maxwidth{0.5\textwidth}]{img1.jpg}
%}
%\vspace{8mm}
%}
%\end{center}

\lipsum

\scieloImageContainerTwoCol{img2.jpg}{\textbf{Figure 1:} A specimen of each species was embedded in paraffin in its entirety, longitudinally sectioned}

\lipsum
\lipsum

\scieloImageContainerTwoCol{img2.jpg}{\textbf{Figure 1:} A specimen of each species was embedded in paraffin in its entirety, longitudinally sectioned}

\lipsum
\lipsum
\lipsum[1]

\begin{scieloReferencesContainer}[REFERÊNCIAS]



\scieloReferencesItem{Jácomo RH, Figueiredo-Pontes LL, Rego EM. Do paradigma molecular ao
impacto no prognóstico: uma visão da leucemia promielocítica aguda.
Rev Assoc Med Bras. 2008;54(1):82-9. DOI: http://dx.doi.org/10.1590/
S0104-42302008000100026}

\scieloReferencesItem{Pinheiro RF, Pelloso LAF, Yamamoto M, Chauffaille MLLF, Bordim JO.
Síndrome Atra: experiência de 10 anos. Rev Bras Canc. 2003;49(1):27-31.}

\scieloReferencesItem{Castaigne S, Chomienne C, Daniel MT, Ballerini P, Berger R, Fenaux P, et al.
All-trans retinoic acid as a differentiation therapy for acute promyelocytic
leukemia. I. Clinical results. Blood. 1990;76(9):1704-9. PMID: 2224119}

\scieloReferencesItem{Rego EM, Santana-Lemos BAA, Tamarozzi MB. Differentiation syndrome
in acute promyelocytic leukemia: pathogenesis and risk factors. Rev Bras
Hematol Hemoter. 2008;30(suppl.2):33-6. DOI: http://dx.doi.org/10.1590/
S1516-84842008000800009}

\scieloReferencesItem{Martín del Pozo M, Cisneros de La Fuente E, Solano F, Martín ML, De la Serna
J. El síndrome del ATRA como complicación del tratamiento en la leucemia
promielocítica aguda. An Med Interna (Madrid). 2001;18(4):195-200.}

\scieloReferencesItem{Franchi L, Baccetti T, De Toffol L, Polimeni A, Cozza P. Phases of the dentition for the assessment of skeletal maturity: a diagnostic performance
study. Am J Orthod Dentofacial Orthop. 2008 Mar;133(3):395-400. PMid: 18331939.}

\scieloReferencesItem{Bullock TH and Horridge GA (1965) Structure
and Function in the Nervous Systems of Invertebrates.
Vol. 2. W.H. Freeman and Co, San
Francisco, 1719 pp.}

\scieloReferencesItem{Carranza S, Baguñà J and Riutort M (1997)
Are the Platyhelminthes a monophyletic primitive
group? An assessment using 18S rDNA
sequences. Mol Biol Evol 14:485-497.}

\scieloReferencesItem{Cebrià F (2008) Organization of the nervous system in the model
planarian Schmidtea mediterranea: An immunocytochemical study. Neurosci Res 61:375-384.}

\scieloReferencesItem{Chien P and Koopowitz H (1972) The ultrastructure of neuromuscular systems in Notoplana acticola, a free-living polyclad flatworm. Z Zellforsch Mikroskop Anat 133:277-288.}


\scieloReferencesItem{Chien PK and Koopowitz H (1977) Ultrastructure of nerve plexus
in flatworms. III. The infra-epithelial nervous system. Cell
Tissue Res 176:335-347.}

\scieloReferencesItem{Day TA, Maule AG, Shaw C and Pax RA (1997) Structure-activity relationships of FMRFamide-related peptides contracting Schistosoma mansoni muscle. Peptides 18:917-921.}

\scieloReferencesItem{Egger B, Gschwentner R and Rieger R (2007) Free-living
flatworms under the knife: Past and present. Dev Genes Evol
217:89-104.}

\scieloReferencesItem{Egger B, Lapraz F, Tomiczek B, Müller S, Dessimoz C, Girstmair
J, Skunca N, Rawlinson KA, Cameron CB, Beli E et al.
(2015) A transcriptomic-phylogenomic analysis of the evolutionary relationships of flatworms. Curr Biol 25:1-7.}

\scieloReferencesItem{Ehlers U (1985) Das Phylogenetische System der Plathelminthes.
Gustav Fischer Verlag, Stuttgart, 317 pp.}

\scieloReferencesItem{Fernandes MC, Alvares EP, Gama P and Silveira M (2003) Serotonin in the nervous system of the head region of the land planarian Bipalium kewense. Tissue Cell 35:479-486.}

\scieloReferencesItem{Forest DL and Lindsay SM (2008) Observations of serotonin and
FMRFamide-like immunoreactivity in palp sensory structures and the anterior nervous system of spionid polychaetes. J Morphol 269:544-551.}

\scieloReferencesItem{Girstmair J, Schnegg R, Telford MJ and Egger B (2014) Cellular
dynamics during regeneration of the flatworm Monocelis sp. (Proseriata, Platyhelminthes). Evo Devo 5:e37.}

\scieloReferencesItem{Golding DM (1992) Polychaeta: Nervous system. In: Harrison
FW and Gardiner SL (eds) Microscopic Anatomy of Invertebrates. Vol. 7. Wiley-Liss, New York, pp 155-179.}

\scieloReferencesItem{Gustafsson MKS, Halton DW, Kreshchenko ND, Movsessian SO,
Raikova OI, Reuter M and Terenina NB (2002) Neuropeptides in flatworms. Peptides 23:2053-2061.}

\scieloReferencesItem{Hadenfeldt D (1929) Das Nervensystem von Stylochoplana
maculata und Notoplana atomata. Z wiss Zool 133:586638.}

\scieloReferencesItem{Halton DW and Gustafsson MKS (1996) Functional morphology
of the platyhelminth nervous system. Parasitology
113:S47-S72.}

\end{scieloReferencesContainer}

\scieloLicenseContainer{License information: This is an open-access article distributed under the terms of the Creative Commons Attribution License, which permits unrestricted use, distribution, and reproduction in any medium, provided the original work is properly cited.}


\end{multicols}
